\documentclass[11pt]{article}
\usepackage[margin=1in]{geometry}
\usepackage{amsmath,amssymb,amsthm,mathtools}
\usepackage[T1]{fontenc}
\usepackage{lmodern,microtype,xurl}
\usepackage[colorlinks=true,linkcolor=blue,urlcolor=blue]{hyperref}
\newtheorem{theorem}{Theorem}[section]
\newtheorem{lemma}[theorem]{Lemma}
\newtheorem{proposition}[theorem]{Proposition}
\theoremstyle{remark}\newtheorem{remark}[theorem]{Remark}
\newcommand{\E}{\mathbb E}
\newcommand{\Prb}{\mathbb P}
\newcommand{\R}{\mathbb R}
\newcommand{\norm}[1]{\left\lVert #1\right\rVert}
\newcommand{\sig}{\mathsf s}
\newcommand{\poly}{\operatorname{poly}}
\title{Shallow targets under deep logistic posteriors}
\author{ChatGPT-6 Astra \and Kaarel H\"anni \and Dmitry Vaintrob}
\date{25 October 2026}
\begin{document}
\maketitle
\begin{abstract}
Consider the standard logistic activation, augmented Boolean inputs, no
internal biases, and a linear readout. We study the actual Gaussian
posterior conditioned on absolute output margin greater than $1/2$ on the
entire cube and exact training labels. Every bounded-weight one-hidden-layer
target has polynomial sample complexity at every finite learner width and
depth in the strong-prior limit. The same holds at every fixed variance
exponent $\gamma>1$. A complementary-pair construction transmits the shallow
computation at energy $16L+O(\log T+1)$ once $L-1\ge\log T$, matching the leading nonconstant
energy lower bound. Centering realized activations and subtracting typical
Gaussian matrices then reduces the finite-scale posterior to a common
class of small-product Lipschitz networks. We prove the required finite
class bound by elementary layer compression. The theorem does not cover
the critical exponent $\gamma=1$ or arbitrary deep teachers, and does not
settle the original tanh question.
\end{abstract}

\section{Model and statement}

Let $X_d=\{-1,1\}^d$, $d\ge1$, and $\sig(a)=(1+e^{-a})^{-1}$. A learner
has $L\ge2$ matrices, $L-1$ hidden layers of common width $N$, input
$(1,x)$, and a linear scalar output $z_W(x)$. There are no internal
biases. The independent Gaussian entry variances are
\[
 t^2/d,\qquad t^2/N,\qquad t^2/N^2
\]
for the first matrix, internal matrices, and readout. We allow a prescribed
architecture-dependent scale $t=t(d,L,N)$ with $0<t\le1$; no constancy
across widths is assumed. For each fixed $\zeta>0$, the finite-scale theorem covers every choice
$t(d,L,N)\le N^{-\zeta/2}$, including
$t(N)=N^{(1-\gamma)/2}$ for each fixed $\gamma=1+\zeta>1$.
For a target $f:X_d\to\{-1,1\}$ and sample $S=((x_i,f(x_i)))_{i=1}^n$,
write
\[
 E_S=\{|z_W(x)|>1/2\text{ for every }x\in X_d\}
       \cap\{\operatorname{sign}z_W(x_i)=f(x_i)\text{ for every }i\},
 \qquad Q_S=\mu(\,\cdot\mid E_S).
\]
Thus both conditioning operations are retained. For iid uniform cube
examples define
\[
 e(S)=2^{-d}\#\{x:Q_S(\operatorname{sign}z_W(x)=f(x))\le2/3\}.
\]
A shallow teacher has one hidden layer of width $C$, all physical weights
bounded by $B\ge1$, and $f(x)z_*(x)\ge1$ everywhere. A fixed positive
teacher margin can instead be normalized by rescaling its readout.

\begin{theorem}\label{thm:main}
There is a universal polynomial $P_0$ with the following property. Every
shallow teacher as above is learned, at every $N\ge C$ and $L\ge2$, by
every limiting sign posterior as $t\downarrow0$, with
\[
 n\ge P_0(d,C,B,L,\varepsilon^{-1},\log(2/\delta))
 \quad\Longrightarrow\quad
 \Prb_S\{e(S)>\varepsilon\}\le\delta.
\]
Here a limiting sign posterior is any subsequential limit on the finite
set of Boolean truth tables, for each fixed finite architecture.
Furthermore, for every fixed $\zeta>0$ there is a polynomial $P_\zeta$
giving the same conclusion for the actual finite-scale posterior,
simultaneously for all
\[
 N\ge C,\quad L\ge2,\quad 0<t\le N^{-\zeta/2}.
\]
Equivalently, the common variance in the isotropic coordinates below is
at most $N^{-(1+\zeta)}$. Polynomial degrees may depend on $\zeta$.
\end{theorem}

The theorem includes both-label samples. Its proof does not use the prior
mass of constant-sign truth tables as a denominator. Constant targets are
included, but are not the reason for the upper bound.

Put
\[
 A_1=\sqrt{d/N}\,W_1,\qquad A_j=W_j\ (2\le j<L),\qquad
 A_L=\sqrt N\,v,\qquad
 \kappa(A)=\sum_{j=1}^L\norm{A_j}_F^2.
\]
All parameter coordinates now have variance $\sigma^2=t^2/N$. The input is
$u_0=(1,x)/\sqrt d$, the normalized hidden states are $u_j=h_j/\sqrt N$,
and the scalar hidden activation is
$\sig(\sqrt N a)/\sqrt N$. It is $1/4$-Lipschitz, and
$\norm{u_j}\le1$ for every hidden layer.

\section{A width-independent finite class bound}

We need only a bound polynomial in depth; no circuit-compression theorem
is used. The activations in this section may differ between neurons.

\begin{lemma}\label{lem:cover}
Fix a finite set $X\subset\R^D$ of radius at most two, with
$\log(2|X|)\le d+1$, and an integer depth $L$. Consider networks whose
hidden scalar activations fix zero and are $1$-Lipschitz, whose product of
matrix Frobenius norms is at most $P\ge1$, and whose final additive scalar
offset has magnitude at most $2P$. The sign functions with absolute
margin at least $3/8$ on $X$ belong to a finite class $\mathcal F$ such that
\[
 \log|\mathcal F|\le\poly(D,d,L,P).
\]
The class is common to all original widths, activation maps, and real
parameter values.
\end{lemma}

\begin{proof}
We first compress one hidden layer. Let $W,V$ be its incoming and outgoing
matrices; let $w_i$ be row $i$ of $W$, $v_i$ column $i$ of $V$, and set
$a_i=\norm{w_i}$, $b_i=\norm{v_i}$,
$A=\sum_i a_i b_i\le\norm W_F\norm V_F$. Omit zero terms. Replacing
$\phi_i$ by $\psi_i(s)=\phi_i(a_i s)/a_i$ and the incoming row by
$w_i/a_i$ gives
\[
 V\phi(Wu)=\sum_i a_i v_i\psi_i((w_i/a_i)u).
\]
Sample $m$ independent indices with probabilities $a_i b_i/A$. For each
sample use incoming row $w_i/a_i$, activation $\psi_i$, and outgoing column
$A v_i/(m b_i)$. Their incoming and outgoing Frobenius norms are exactly
$\sqrt m$ and $A/\sqrt m$, so the product does not increase. The output
is an unbiased average of vectors of norm at most $A\norm u$.

If $S$ is the product of preceding norms, the current states on $X$ have
norm at most $2S$. The expected norm error of the average is at most
$2AS/\sqrt m$. Bounded-difference concentration, followed by a union bound
over $X$, gives a choice of the sampled indices with preactivation error
at most $C AS\sqrt{(d+1)/m}$ simultaneously on $X$. The suffix is
Lipschitz with constant at most the product of its matrix norms. Hence
the final scalar error is at most $C P\sqrt{(d+1)/m}$, irrespective of
the original layer widths. This also holds when $V$ is the final readout.
Compress every hidden layer in succession, using
\[
 m=\left\lceil C P^2L^2(d+1)\right\rceil.
\]
The total output error is at most $1/16$, after increasing $C$, and every
hidden width is at most $m$. Every compression preserves zero-fixing
$1$-Lipschitz activations and does not increase the global product.

Now normalize each hidden matrix to Frobenius norm one. If its original
norms are $a_1,\ldots,a_{L-1}$, divide the hidden state at layer $j$ by
$S_j=\prod_{i\le j}a_i$, replace its scalar activation by
$\phi(S_j s)/S_j$, and multiply the final readout by $S_{L-1}$.
All activations remain zero-fixing and $1$-Lipschitz; the readout norm is
at most $P$, and all hidden states have norm at most two. If a hidden
matrix has norm zero, the network is constant and can be listed separately.

Round hidden matrices in Frobenius norm to accuracy
$\eta=c/(PL)$, and round the readout and final offset to accuracy $c$.
Store every scalar activation on a uniform grid of spacing
$\eta'=c/(PL\sqrt m)$ in $[-7,7]$, rounding its values to this accuracy
and using linear interpolation. The table differs from the original
activation by at most $2\eta'$ on that interval. The elementary state
error recurrence, using the original activation's Lipschitz constant one,
shows that all rounded states have norm below three and the final output
changes by at most $1/16$, for a sufficiently small universal $c$.
For example the accumulated hidden error is bounded by
$C L(3\eta+\sqrt m\eta')$, since the rounded matrix norms are at most
$1+\eta$. Readout and offset errors can be made smaller fixed constants
by reducing $c$. All preactivations stay in the stored interval.

There are polynomially many matrix entries and table entries in
$D,d,L,P$, each with polynomial logarithmic range and precision. Thus
the total description length is polynomial in these parameters. The
two approximations change every output by less than $1/8$, preserving
the signs of the stated margin class. Listing all descriptions proves
the lemma; no precision hypothesis on the original activations is needed.
\end{proof}

\section{A matching-energy complementary-pair wire}

A nonconstant globally margined classifier has opposite labels on some
cube edge. The normalized inputs at its endpoints have distance
$2/\sqrt d$, while the outputs differ by more than one. Since logistic
has derivative at most $1/4$, AM--GM gives the deterministic bound
\begin{equation}\label{eq:lower}
 \prod_{j=1}^L\norm{A_j}_F>4^{L-1}\sqrt d/2,\qquad
 \kappa(A)\ge16L\exp\bigl((\log d-\log64)/L\bigr).
\end{equation}
In particular the leading energy coefficient is $16$. This is not a
finite-temperature posterior rate calculation.

\begin{lemma}\label{lem:scalar}
For $D>0$ and $m\ge2$, put $D'=\max(D,4)$ and
$q=(D'/4)^{1/(2m)}$. There are positive weights
$\alpha,b_2,\ldots,b_{m-1},c$ such that
\[
 r_1=\tanh\alpha,\quad r_s=\tanh(b_s r_{s-1}),\quad
 cr_{m-1}=1/2,
\]
and
\[
 D\alpha^2+\sum_{s=2}^{m-1}b_s^2+c^2\le(m+12)q^2,
 \qquad c=q.
\]
Moreover every consecutive product of the $b_s$ is at most
$Cm\sqrt{D'/4}$.
\end{lemma}

\begin{proof}
Set $k_s=\min(s,m-s)$ and $r_s=q^s/(\sqrt{D'}k_s)$. Then
$r_s\le1/(2q k_s)\le1/2$ and $r_{m-1}=1/(2q)$. Take
$\alpha=\operatorname{atanh}r_1$, $b_s=\operatorname{atanh}(r_s)/r_{s-1}$,
and $c=q$. For $0<u\le1/2$, the power series gives
$(\operatorname{atanh}(u)/u)^2\le1+u^2$.
The reciprocal pairs in $R_s=k_{s-1}/k_s$ obey
\[
 ((k+1)/k)^2+(k/(k+1))^2-2\le4/k^2.
\]
Consequently $\sum R_s^2\le m-2+7$, $R_s\le2$, and
$\sum k_s^{-2}<4$. These imply
\[
 D\alpha^2\le q^2+1/4,\qquad
 \sum b_s^2\le q^2(m-2+7)+4,\qquad c^2=q^2,
\]
whose sum is at most $(m+12)q^2$. Empty sums when $m=2$ cause no change.
Finally
\[
 b_s=q\frac{k_{s-1}}{k_s}\frac{\operatorname{atanh}r_s}{r_s},
 \qquad \sum_s r_s^2\le1.
\]
The $k$-ratios telescope, $q^m=\sqrt{D'/4}$, and the remaining factors
have product at most $\exp(\frac12\sum r_s^2)$. This proves the partial
product bound.
\end{proof}

Let $g(x)=u\cdot h_1(x)$ be a shallow witness with signed margin one,
first energy $E_1$, and weighted readout energy $D=N\norm u^2$.
For $N\ge2$, use $M=2\lfloor N/2\rfloor$ active neurons in each added
hidden layer, split into equal groups. Let $a$ have entries $2/M$ and
$-2/M$ on the groups and zero on any surplus coordinate. Then
$\norm a^2=4/M$, and
\[
 \sig(2z)-\sig(-2z)=\tanh z.
\]
Apply Lemma~\ref{lem:scalar} with first coefficient $D/4$ and $m=L-1$.
The first added layer has rows $\pm2\alpha u$ and energy at most
$4D\alpha^2$. Each later added layer has rows $\pm2b_s a^T$ and energy
$16b_s^2$. The final readout $2ca$ has weighted energy
$16(N/M)c^2\le32c^2$. Monotonicity of the scalar chain gives global
signed output at least one. Thus
\begin{equation}\label{eq:wire}
 \kappa(A_*)\le E_1+16(L+12)\max(1,D/16)^{1/(L-1)}.
\end{equation}
Odd widths have only an endpoint cost; they do not increase the per-layer
coefficient $16$.

Choose an integer
\[
 M_0\ge\max\{C,2,d+1,dC(d+1)B^2\}
\]
of polynomial size. Zero-padding the teacher to width $M_0$ makes its
first energy at most one. Balanced cloning to any $N\ge M_0$ gives
$E_1\le2$ and $D\le D_0:=2M_0CB^2$: each replication count lies between
$N/(2M_0)$ and $2N/M_0$. Put $T=\max(1,D_0/16)$. The construction gives
\begin{equation}\label{eq:K}
 K_L=2+16(L+12)T^{1/(L-1)},\qquad \kappa(A_*)\le K_L.
\end{equation}
In particular, for every $L\ge3$,
\begin{equation}\label{eq:P}
 4\left(\frac{K_L+1}{16L}\right)^{L/2}
 \le4e^{13/2}T^{3/4}.
\end{equation}
Indeed $(K_L+1)/16\le(L+13)T^{1/(L-1)}$ and
$L/(2(L-1))\le3/4$. For $L=2$ the cloned shallow witness directly has
energy at most $2+D_0$, which gives a polynomial bound as well.

\begin{lemma}\label{lem:stable}
The above witness can be chosen so that perturbing every parameter block
in operator norm by at most $\eta$ changes its output on the entire cube
by at most $C L^2T\eta$. The same bound, with an enlarged constant,
holds for the direct $L=2$ witness.
\end{lemma}

\begin{proof}
Use the upper bound $D_0$ in the scalar lemma. The first added layer's
normalized Lipschitz gain $\norm{A_2^*}_{\rm op}/4$ is at most $Cq$;
the last readout norm is at most $Cq$. The first teacher matrix has norm
at most $\sqrt2$. The intervening gains are exactly $b_s$ in operator
norm, and Lemma~\ref{lem:scalar} bounds every partial product by
$CL\sqrt T$. Including the endpoint factors gives a bound $CLT$ for
every suffix affecting an injected error.

Normalized logistic hidden states of both the witness and the perturbed
network have norm at most one. Therefore their state differences satisfy
\[
 e_j\le(\norm{A_j^*}_{\rm op}/4)e_{j-1}+\eta/4,
\]
with first forcing at most $\sqrt2\eta/4$. Summing the $L$ forced errors,
and treating the final linear readout similarly, proves the result.
The proof uses the explicit partial products; total energy alone would
not supply this polynomial stability bound.
\end{proof}

\section{Centering and finite-scale Gaussian subtraction}

Fix an anchor $x_0\in X_d$. Center every normalized hidden state at that
anchor. Writing $b_{j,i}=(A_j u_{j-1}(x_0))_i$, the centered scalar maps
are
\[
 a_{b,N}(r)=\frac{\sig(\sqrt N(b+4r))-\sig(\sqrt N b)}{\sqrt N}.
\]
They fix zero and are $1$-Lipschitz. The centered linear blocks are
$A_j/4$ for hidden layers and $A_L$ for the readout. Thus the centered
product is at most
\begin{equation}\label{eq:centerproduct}
 P(A)=4^{-(L-1)}\prod_j\norm{A_j}_F
 \le4\left(\frac{\kappa(A)}{16L}\right)^{L/2}.
\end{equation}
The centered input $(0,x-x_0)/\sqrt d$ has norm at most two.
The output is $z_A(x_0)$ plus this centered computation. If the truth
table is nonconstant, an opposite-sign point implies
\begin{equation}\label{eq:offset}
 |z_A(x_0)|\le2P(A).
\end{equation}
Constants can be listed separately.

For $N\ge M_0$, let $\mathcal T$ be the intersection of the operator-norm
balls of radius $tb$ in all parameter blocks, where
\[
 b^2=16\log9+8\log(8L).
\]
This is a closed convex centrally symmetric set with $\mu(\mathcal T)
\ge1/2$. For completeness, a $1/4$-net in each unit sphere has cardinality
at most $9^r$ in dimension $r$. The bilinear Gaussian tail and the net
bound give
$\Prb(\norm G_{\rm op}>tb)\le2\,9^{r+c}\exp(-Nb^2/8)$ for an
$r\times c$ block of variance $t^2/N$. Here $r+c\le2N$, since
$N\ge d+1$, so a union bound over the $L$ blocks proves the claim.

Assume
\begin{equation}\label{eq:smallbulk}
 tb\le\xi:=c/(L^2T)
\end{equation}
for a sufficiently small universal $c$. Lemma~\ref{lem:stable} puts
$A_*+\mathcal T$ inside global signed margin $3/4$, hence inside every
$E_S$ labeled by the target. Gaussian translation and symmetry give
\begin{equation}\label{eq:denom}
 \mu(E_S)\ge\mu(A_*+\mathcal T)
 \ge\tfrac12\exp[-K_L/(2\sigma^2)].
\end{equation}
Indeed the translated density is the original density multiplied by
$\exp[-\norm{A_*}^2/(2\sigma^2)-\langle A_*,A\rangle/\sigma^2]$;
symmetry replaces the last factor by a hyperbolic cosine of mean at least
one. Gaussian isoperimetry for a set of mass at least $1/2$ gives
$\mu\{\operatorname{dist}(A,\mathcal T)>r\}\le\exp[-r^2/(2\sigma^2)]$.
Dividing by \eqref{eq:denom} proves the actual posterior estimate
\begin{equation}\label{eq:distance}
 Q_S\{\operatorname{dist}(A,\mathcal T)^2>K_L+1\}
 \le2\exp[-1/(2\sigma^2)].
\end{equation}
Only this standard Gaussian isoperimetric inequality is used as a
probabilistic input beyond elementary concentration.

On the complementary event, project $A$ onto $\mathcal T$ and write
$A=U+B$, where $\kappa(U)\le K_L+1$ and every block of $B$ has operator
norm at most $tb$. Center the \emph{realized} network $A$ at $x_0$ and
freeze its exact maps $a_{b,N}$. Replace its centered linear matrices by
those of $U$, keeping both the frozen maps and the actual anchor output.
This surrogate need not be a native logistic network. It belongs to the
larger Lipschitz class of Lemma~\ref{lem:cover}, which is all that counting
requires. The anchor maps are not recomputed after subtraction.

To bound the change, divide all $U$ blocks, including the readout, by four
and place a factor four outside the final output. Set
$K'=(K_L+1)/16$ and $\eta=tb/4$. A telescoping replacement of linear
blocks, using zero-preserving $1$-Lipschitz maps and input radius two,
gives error at most
\begin{equation}\label{eq:subtract}
 8L\eta
 \left(\frac{(\sqrt{K'}+\sqrt L\eta)^2}{L-1}\right)^{(L-1)/2}.
\end{equation}
To verify it, each of the $L$ terms is bounded by $8\eta$ times the
product of the other $L-1$ block norms, each enlarged by at most $\eta$.
Their squared sum is at most $(\sqrt{K'}+\sqrt L\eta)^2$; apply AM--GM.
For $L\ge3$, \eqref{eq:smallbulk} and
$K'\le(L+13)T^{1/(L-1)}$ bound the displayed power by $C\sqrt T$.
Indeed the added energy is bounded by a universal constant, after which
the ratio is at most $(1+C/(L-1))T^{1/(L-1)}$.
Thus \eqref{eq:subtract} is at most $1/8$, after decreasing $c$.
For $L=2$, use its direct bound $K'\le C(1+T)$ and the same conservative
condition \eqref{eq:smallbulk}.

The surrogate preserves every sign of the globally margined posterior
network and has margin greater than $3/8$. Its product is polynomial in
$T$ by \eqref{eq:P}. Since every frozen map fixes zero, its centered
computation vanishes at the anchor. If its truth table is nonconstant,
the retained offset is automatically at most twice its own product,
by the same opposite-sign argument as \eqref{eq:offset}. Lemma~\ref{lem:cover}
therefore supplies one class $\mathcal F$ with
\begin{equation}\label{eq:commonclass}
 \log|\mathcal F|\le\poly(d,C,B,L),\qquad
 Q_S\{\operatorname{sign}z_A\notin\mathcal F\}
 \le2e^{-1/(2\sigma^2)}.
\end{equation}
The class is independent of the sample, width, and realized anchor values.

\section{Small widths, cold limits, and statistical conversion}

We record the elementary finite-width estimate needed below the
large-width threshold. It also avoids any appeal to sigmoid VC theory.

\begin{lemma}\label{lem:finitewidth}
Suppose a width-$N$, depth-$L$ logistic target witness has signed margin
one and isotropic energy at most $K_0\ge1$. For $0<t\le1$, its posterior
has sufficient sample size polynomial in
\[
 d,N,L,K_0,\varepsilon^{-1},\log(2/\delta),
\]
uniformly in $t$.
\end{lemma}

\begin{proof}
The parameter dimension is $p=N(d+2)+(L-2)N^2$. On an energy ball
$\kappa\le R$, the map from parameters to any cube output has a
Lipschitz constant $\Lambda_R$ with
\[
 \log(2+\Lambda_R)\le CL\log(2+R+L).
\]
For example differentiate layer by layer: all hidden states have norm
at most one, the input norm is at most $\sqrt2$, each derivative is
bounded by $1/4$, and every downstream block norm is at most $\sqrt R$.
This proves a bound $C L\max(1,\sqrt R)^L$, which suffices.

Let $K=2K_0+2$ and $r=\min(1,1/(4\Lambda_K))$. The ball of radius $r$
around the witness has signed margin at least $3/4$ and energy at most
$K$. Its inscribed cube with half-side $\min(r/\sqrt p,\sigma)$ gives
\[
 \mu(E_S)\ge\exp[-K/(2\sigma^2)-pT_0],\qquad
 T_0=\log\!\left(C\max\{1,\sqrt{p/N}/r\}\right).
\]
Here the clipping at $\sigma$ is essential: since $\sigma\le1/\sqrt N$,
the volume-to-Gaussian-scale ratio introduces no $\log(1/t)$ term.
The isotropic Gaussian tail
$\mu\{\kappa>R\}\le2^{p/2}e^{-R/(4\sigma^2)}$ shows that the posterior
mass outside
\[
 R_\beta=2K+\frac4N\{p(T_0+\tfrac12\log2)+\log(1/\beta)\}
\]
is at most $\beta$. This radius is polynomial in the displayed resources
and $\log(1/\beta)$.

A parameter net of mesh $1/(8\Lambda_{R_\beta+1})$ in that ball has
log-cardinality polynomial in the same resources. Every globally margined
network in the ball has exactly the signs of a net representative, since
the output approximation error is below $1/4$. Thus a common finite
class contains posterior mass at least $1-\beta$. The finite-class
argument given below completes the proof.
\end{proof}

For $C\le N<M_0$ and $N\ge2$, zero-padding the shallow teacher to width
$N$ gives first energy at most $dC(d+1)B^2/N$ and readout energy at most
$NCB^2$. For $L\ge3$, applying \eqref{eq:wire} gives a witness energy
polynomial in $d,C,B,L,N$; for $L=2$, use the direct shallow witness.
Hence Lemma~\ref{lem:finitewidth}
handles these widths. For $N=1$ every logistic truth table is constant,
since the last hidden value is positive and the readout is scalar;
one example fixes its sign everywhere.

For the finite-scale part of Theorem~\ref{thm:main}, fix $\zeta>0$.
There is a threshold $N_*\ge M_0$, polynomial in
$d,C,B,L,\log(2/\beta)$ with degree depending only on $\zeta$, such that
every $N\ge N_*$ satisfies both \eqref{eq:smallbulk} and
$2e^{-1/(2\sigma^2)}\le\beta$ whenever $t\le N^{-\zeta/2}$.
This follows from $b=O(\sqrt{\log(2L)})$, $T=\poly(d,C,B)$, and
$\sigma^2\le N^{-(1+\zeta)}$. Equation~\eqref{eq:commonclass} handles
all widths above $N_*$. For widths below it, the just-described witness
and Lemma~\ref{lem:finitewidth} give a polynomial bound because $N_*$
itself is polynomial. Take the union of the large-width class with the
finite parameter-net classes for the integers $C\le N<N_*$. Its
log-cardinality increases by at most $\log(N_*+1)$ beyond the largest
individual bound, and remains polynomial. This gives one sample event
for all widths and eligible scales at the fixed depth.

For the cold statement, at every fixed finite architecture a strict-margin
witness of energy at most $K$ implies
\begin{equation}\label{eq:cold}
 Q_S\{\kappa>K+\rho\}\longrightarrow0\qquad(\sigma\downarrow0)
\end{equation}
for each $\rho>0$, uniformly over target-labeled samples. A sufficiently
small fixed neighborhood of the witness retains its signs and margin
and has energy at most $K+\rho/2$. Its Gaussian mass is bounded below by
a constant times $\sigma^{-p}e^{-(K+\rho/2)/(2\sigma^2)}$.
The Gaussian energy tail above $K+\rho$ is at most a polynomial in
$\sigma^{-1}$ times $e^{-(K+\rho)/(2\sigma^2)}$, so their ratio tends
to zero. With \eqref{eq:K}, \eqref{eq:centerproduct}, and
Lemma~\ref{lem:cover}, every limiting sign posterior at $N\ge M_0$ is
supported on a common class of polynomial log-cardinality. The preceding
finite-width argument covers the remaining widths uniformly in the
cold limit. Passing Gibbs error through a sign-posterior limit is valid
on the finite truth-table simplex; no continuity of the inclusive
confidence threshold is assumed.

Finally, let a sample-independent class of size $M$ contain posterior
mass at least $1-\beta$. With probability at least $1-\delta$, every
sample-consistent member of this class has true error at most $\alpha$,
provided
\[
 n\ge\frac{\log M+\log(1/\delta)}{\alpha}.
\]
Every draw from $Q_S$ is sample consistent. Its expected true error is
therefore at most $\alpha+\beta$. If correct-label confidence is at most
$2/3$, its wrong-label probability is at least $1/3$, so
$e(S)\le3(\alpha+\beta)$. Taking $\alpha=\beta=\varepsilon/12$
leaves strict slack at the stated threshold. For a limiting posterior
supported on the class take $\beta=0$.

To obtain the simultaneous assertion across depths, fix the other target
and statistical resources and put $\delta_L=\delta/[L(L-1)]$ for
$L\ge2$. At each depth the argument above already gives one event for
all widths and eligible scales. For a fixed sample size $n$, union bound
over the depths whose sufficient sample condition is met, using
$\delta_L$ in place of $\delta$. Their failure probabilities sum to at
most $\delta$. Since
$\log(2/\delta_L)=\log(2/\delta)+\log(L(L-1))$, enlarging the
existential polynomial sample budget absorbs this additional term.
The same reasoning applies to the cold classes. This proves the stated
simultaneity without asserting a depth-uniform positive cold threshold,
and completes Theorem~\ref{thm:main}.

\section{Scope and a concrete shallow example}

Parity is included in the shallow-target hypothesis. To see this
directly, put $k(x)=(d+\sum_i x_i)/2\in\{0,\ldots,d\}$ and
$y(k)=(-1)^k$. This is full parity up to the fixed sign $(-1)^d$,
which can be absorbed into the readout. The step expansion
\[
 y(k)=y(0)+\sum_{j=1}^d[y(j)-y(j-1)]\,\mathbf1_{\{k>j-1/2\}}
\]
is approximated on these integer values by replacing each step with
$\sig(A(k-j+1/2))$, where $A=2\log(16d)$.
The total approximation error is at most $2d e^{-A/2}\le1/8$.
The constant term is realized by a zero-weight hidden unit with output
$1/2$. Multiplying the final output by two gives signed margin at least
one. There are $d+1$ hidden units and weights bounded by $O(d\log(2d))$.

Thus neither parity nor another bounded-weight shallow logistic target
can yield an exponential-in-depth sample lower bound in the regimes
of Theorem~\ref{thm:main}. The proof leaves the exact critical exponent
$\gamma=1$, and a fixed positive canonical prefactor $t$, open: then
the typical Gaussian block size $tb$ need not satisfy
\eqref{eq:smallbulk}. It also leaves arbitrary deep teachers open.
The offset of standard logistic is handled here by centering and
complementary transmission; this result does not assert the analogous
theorem for other activations or resolve the original tanh problem.

\begin{thebibliography}{9}
\bibitem{borell}
C. Borell, \emph{The Brunn--Minkowski inequality in Gauss space},
Inventiones Mathematicae \textbf{30} (1975), 207--216.
\end{thebibliography}
\end{document}
